\chapter{Anexos del Subdocumento 12}

Los 230 paquetes de la EDT tienen responsable en el Formulario T-14. El Anexo 12.A cruza esos paquetes con las horas propias de los titulares del Formulario T-15; los Anexos 12.B y 12.C identifican a cada persona y su dedicación; el Anexo 12.D registra el compromiso de cada integrante con su nominación del Formulario T-8. Los meses se cuentan desde el inicio contractual, no desde el ingreso de cada integrante.

\section{Anexo 12.A — Matriz de asignación de responsabilidades}

La Tabla~\ref{tab:12A-matriz} se lee por producto. A indica que el rol responde por al menos un paquete del T-14; E, que ejecuta trabajo con horas propias del T-15, sin responder por paquetes de ese producto. Si concurren ambas condiciones, prevalece A. Una celda vacía indica que no existe ninguna de las dos asignaciones.

Las columnas siguen el orden del T-8: JP, Jefe de Proyecto; ARQ, Arquitecto de Solución; SEG, Encargado de Seguridad de la Información; DAT, Líder de Datos; DES, Líder de Desarrollo; CAL, Líder de Calidad y Pruebas; SRE, Líder de Operación / SRE; IMP, Líder de Implantación y Gestión del Cambio; LF, Líder Funcional, y LINT, Líder de Integración.

\begin{anexohorizontal}
\setlength{\textwidth}{\linewidth}
\parte{tablas/tab_12A_matriz}

La última fila cuenta paquetes, no marcas A. Un producto puede tener varios responsables porque agrupa varios paquetes; cada paquete conserva un único responsable. Los totales suman 230. Las marcas E del Líder Funcional y del Líder de Integración identifican trabajo efectivo sin duplicar esa cuenta.
\end{anexohorizontal}

\section{Anexo 12.B — Currículos del equipo}

Los ocho primeros currículos corresponden al equipo clave y los dos últimos a los otros roles mínimos (Bases Técnicas Transversales, sección 19.2, pp.~33--34). La experiencia se vincula con las referencias del Formulario T-6: el sistema híbrido de reparto de LogiNacional S.A., la plataforma de telemetría y cadena de frío de FarmaRed S.A. y el sistema móvil de preventa y ruteo de Comercializadora Lácteos del Sur S.A. Los años de obtención y las vigencias se indican por credencial; no representan fechas de ingreso al proyecto.

\subsection{Alex Aravena — Jefe de Proyecto}

Ingeniero Civil en Informática, con 15 años de experiencia. Dirigió el Sistema Híbrido de Reparto y Gestión Logística de LogiNacional S.A. (2021--2022), con 420 camiones, operación continua y un SLA de 99,5~\%. También dirigió el Sistema Móvil de Preventa y Ruteo Desconectado de Comercializadora Lácteos del Sur S.A. (2023--2024). Su especialidad es la dirección de contratos con hitos de pago y control por valor ganado.

\textbf{Acreditación.} Project Management Professional (PMP) (PMI): 2019, renovada en 2025. Hasta 2028; ciclo de 3 años con 60 PDU. Certified ScrumMaster (CSM) (Scrum Alliance): 2021, renovada en 2025. Hasta 2027; renovación cada 2 años.

\textbf{Asignación.} Jefe de Proyecto; meses 1--56. Dedicación declarada: 100~\% en implementación y 22~\% en Operación. Conduce la implementación con dedicación exclusiva y representa a LafroX ante el CLIENTE; en Operación mantiene la dirección contractual.

\subsection{Bastián Trejo — Arquitecto de Solución}

Ingeniero Civil en Informática, con 12 años de experiencia. Diseñó la arquitectura híbrida de LogiNacional S.A., con servicios AWS y procesamiento local en 14 centros de distribución. Fue arquitecto de la Plataforma de Telemetría y Cadena de Frío IoT de FarmaRed S.A. Su especialidad es el diseño de sistemas híbridos y su evaluación mediante revisiones Well-Architected.

\textbf{Acreditación.} AWS Certified Solutions Architect -- Professional (Amazon Web Services): marzo de 2025. Hasta marzo de 2028; 3 años.

\textbf{Asignación.} Arquitecto de Solución; meses 1--56. Dedicación media: 39~\% en implementación y 2~\% en Operación. Concentra su trabajo de diseño en los meses 1 a 7 y participa durante todo el contrato en el Comité de Arquitectura mensual. Los porcentajes son promedios individuales del perfil 03S, compartido entre dos arquitectos senior.

\subsection{Álvaro Catalán — Encargado de Seguridad de la Información}

Ingeniero en Computación, con 12 años de experiencia. Tuvo a cargo la seguridad de la plataforma de FarmaRed S.A., con datos regulados y un SLA de 99,9~\%. Participa en la administración del sistema de gestión de seguridad de la información de LafroX, certificado en ISO/IEC 27001. Su práctica combina respuesta a incidentes y continuidad del servicio.

\textbf{Acreditación.} CISSP (ISC2): 2018, ciclo vigente desde 2024. Hasta 2027; ciclo de 3 años con 120 CPE.

\textbf{Asignación.} Encargado de Seguridad de la Información; meses 1--56. Dedicación media: 40~\% en implementación y 5~\% en Operación. Mantiene la supervisión de seguridad durante ambas fases y responde por la aceptación de las pruebas independientes.

\subsection{Leandro Chamorro — Líder de Datos}

Ingeniero Civil Informático, con 10 años de experiencia. Modeló y migró el historial de temperatura y lotes de FarmaRed S.A. y los maestros de clientes, productos y rutas de Comercializadora Lácteos del Sur S.A. Su especialidad es la conciliación de migraciones: comprueba que los datos de destino conserven la información de origen.

\textbf{Acreditación.} AWS Certified Data Engineer -- Associate (Amazon Web Services): abril de 2026. Hasta abril de 2029; 3 años.

\textbf{Asignación.} Líder de Datos; meses 1--20 y apoyo puntual en los meses 21--56. Dedicación media: 54~\% en implementación y 1~\% en Operación. Mantiene participación durante toda la implementación y presta apoyo puntual en Operación.

\subsection{Tomás Pérez — Líder de Desarrollo}

Ingeniero Civil en Informática, con 10 años de experiencia. Lideró el desarrollo móvil en Kotlin para la preventa y el ruteo sin conexión de Comercializadora Lácteos del Sur S.A., y el desarrollo de servicios Laravel (PHP) y aplicaciones Angular en LogiNacional S.A. Su especialidad es la integración y entrega continua de aplicaciones.

\textbf{Acreditación.} AWS Certified Developer -- Associate (Amazon Web Services): septiembre de 2025. Hasta septiembre de 2028; 3 años.

\textbf{Asignación.} Líder de Desarrollo; meses 1--56. Dedicación declarada: 100~\% en implementación y 15~\% en Operación. Dirige la construcción durante toda la implementación y entrega a Operación los antecedentes para desplegar y recuperar el servicio.

\subsection{Maximiliano Miño — Líder de Calidad y Pruebas}

Ingeniero en Informática, con 9 años de experiencia. Lideró el trabajo de calidad en LogiNacional S.A. y Comercializadora Lácteos del Sur S.A. Su experiencia comprende automatización de pruebas, pruebas de carga y aceptación con usuarios operacionales.

\textbf{Acreditación.} ISTQB Certified Tester Foundation Level v4.0 (ISTQB): 2023. Sin vencimiento. ISTQB Certified Tester Advanced Level Test Management (ISTQB): 2025. Sin vencimiento.

\textbf{Asignación.} Líder de Calidad y Pruebas; meses 1--56. Dedicación media: 17~\% en implementación y 1~\% en Operación. Mantiene participación en ambas fases para dirigir las pruebas y respaldar la aceptación con evidencia.

\subsection{Guillermo Castillo — Líder de Operación / SRE}

Ingeniero Civil en Informática, con 10 años de experiencia. Participó en la operación continua de LogiNacional S.A. y FarmaRed S.A., con SLA de 99,5~\% y 99,9~\%, respectivamente. Su especialidad es la confiabilidad de plataformas AWS y la gestión de incidentes, problemas y cambios bajo ISO/IEC 20000-1.

\textbf{Acreditación.} ITIL 4 Foundation e ITIL 4 Specialist: Create, Deliver and Support (PeopleCert): junio de 2025. Hasta junio de 2028; renovación cada 3 años.

\textbf{Asignación.} Líder de Operación / SRE; meses 1--56. Dedicación declarada: 33~\% en implementación y 100~\% en Operación. Participa desde el inicio y actúa como gerente de servicio dedicado durante los meses 21 a 56.

\subsection{Patricio Henríquez — Líder de Implantación y Gestión del Cambio}

Ingeniero Civil Industrial, con 11 años de experiencia. Condujo la implantación en bodegas y rutas de LogiNacional S.A., con 420 camiones en tres regiones, y la implantación de preventa en Comercializadora Lácteos del Sur S.A. Ha capacitado a conductores, preventistas y operadores, y ha trabajado en gestión del cambio con sindicatos.

\textbf{Acreditación.} Prosci Certified Change Practitioner (Prosci): 2022. Sin renovación periódica.

\textbf{Asignación.} Líder de Implantación y Gestión del Cambio; meses 2--4 y 8--20, con apoyo puntual en los meses 21--56. Dedicación media: 51~\% en implementación y 5~\% en Operación. Concentra el acompañamiento a usuarios en las marchas blancas y el cierre de los meses 13 a 22.

\subsection{Carolina Fuentes — Líder Funcional}

Ingeniera Civil Industrial, con 14 años de experiencia. Trabajó en distribución mayorista y consumo masivo. Fue jefa de operaciones de un centro de distribución durante seis años y consultora funcional en LogiNacional S.A., en procesos de cruce de andén y facturación en ruta. Su especialidad es traducir la operación de bodega y reparto en requisitos funcionales.

\textbf{Acreditación.} Certified Supply Chain Professional (CSCP) (ASCM): 2020. Mantenimiento cada 5 años, vigente hasta 2030 (renovada en 2025).

\textbf{Asignación.} Líder Funcional; meses 1--20. Dedicación media: 100~\% en implementación, sin asignación en Operación. Pertenece al grupo de Arquitectura e integración de planta y cubre el análisis funcional durante toda la implementación; no tiene horas asignadas en Operación.

\subsection{Rodrigo Valdés — Líder de Integración}

Ingeniero Civil en Informática, con 11 años de experiencia. Integró el ERP, los documentos tributarios electrónicos y los intercambios EDI y AS2 con cadenas de comercio minorista en LogiNacional S.A. En FarmaRed S.A. integró sensores y sistemas heredados. Su especialidad son las interfaces entre plataformas de distinta generación.

\textbf{Acreditación.} El perfil se acredita por experiencia en EDI, AS2 y ERP heredados; no exige certificación.

\textbf{Asignación.} Líder de Integración; meses 1--20 y apoyo puntual en los meses 21--56. Dedicación media: 37~\% en implementación y 1~\% en Operación. Pertenece al grupo de Arquitectura e integración de planta; cubre la implementación completa y apoya en forma puntual la Operación.

\section{Anexo 12.C — Dedicación por rol y fase}

Una jornada completa equivale a 128 HH programadas por mes (Formulario T-15, sección 4.1). Para un titular, la fracción media de jornada es \(d_{\mathrm{impl}} = \mathrm{HH}_{\mathrm{impl}}/(20\times128)\) en implementación y \(d_{\mathrm{oper}} = \mathrm{HH}_{\mathrm{oper}}/(36\times128)\) en Operación. El perfil 03S reúne dos plazas senior. La carga media individual se estima dividiendo las horas del perfil entre ambas plazas antes de aplicar el denominador de cada fase.

La Tabla~\ref{tab:12C-dedicacion} conserva las HH del Formulario T-15 y muestra la dedicación declarada en la nominación. La Tabla~\ref{tab:12C-jornada} compara el compromiso de jornada completa con las HH que el T-15 valoriza en esa fase para los tres titulares con dedicación completa. El guion indica una fase sin asignación.

\begin{anexohorizontal}
\setlength{\textwidth}{\linewidth}
\parte{tablas/tab_12C_dedicacion}

\begin{tablalafrox}{Jornada completa comprometida y carga media imputada. Fuente: cálculo propio a partir de las horas por perfil y la jornada de 128 HH mensuales del Formulario T-15, sección 4, y las dedicaciones completas del Formulario T-8, sección 1.}{tab:12C-jornada}
  {Y R{2.2cm} R{2.2cm} C{2.65cm} R{2cm}}
  {\cab{Titular y fase} & \cab{Compromiso (HH)} & \cab{Imputación (HH)} & \cab{Carga media imputada} & \cab{Diferencia (HH)}}
  Jefe de Proyecto, implementación & 2.560 & 2.582,10 & 100,86~\% & +22,10 \\
  Líder de Desarrollo, implementación & 2.560 & 2.565,82 & 100,23~\% & +5,82 \\
  Líder de Operación / SRE, Operación & 4.608 & 4.596 & 99,74~\% & −12 \\
\end{tablalafrox}

El compromiso declarado al 100~\% equivale a 128 HH mensuales: 2.560 HH en los meses 1 a 20 y 4.608 HH en los meses 21 a 56. La diferencia se explica porque la tabla de perfiles del T-15 agrupa las HH según la fase de valorización de cada paquete, y algunos paquetes de gestión y cierre cruzan el mes 21, mientras que el calendario mensual asigna las horas al mes calendario en que se ejecutan (Formulario T-15, sección 4). En ambas vistas el total de los 56 meses es el mismo: 3.594,97 HH de Jefatura, 3.257,02 HH de Desarrollo y 5.431,35 HH de SRE. Datos e Integración prestan apoyo puntual en Operación con la dedicación declarada para esa fase.
\end{anexohorizontal}

\section{Anexo 12.D — Compromiso de los integrantes}

Los diez titulares asumen el compromiso que sigue, asociado al rol, la dedicación y el período de su nominación.

\begin{quote}
Me comprometo a desempeñar el rol que consta junto a mi nombre, con la dedicación y durante los meses allí indicados. Resguardaré la confidencialidad de la información del CLIENTE y mantendré disponibles para el equipo autorizado los antecedentes necesarios para continuar mi trabajo. Mi participación no se dará por sustituida sin autorización previa y escrita del CLIENTE. Si se autoriza mi reemplazo, efectuaré el traspaso a una persona de perfil igual o superior, con al menos 15 días hábiles de traslape, sin costo para el CLIENTE.
\end{quote}

La Tabla~\ref{tab:12D-compromisos} identifica a los integrantes que asumen el compromiso. Para el Arquitecto de Solución, el promedio individual incluye la concentración en diseño y la participación permanente en el Comité de Arquitectura. Las condiciones de sustitución se rigen por las Bases Administrativas (arts.~76.1 y 76.2, p.~39).

\parte{tablas/tab_12D_compromisos}

El compromiso abarca tanto a los ocho integrantes clave como a los dos titulares adicionales, sin alterar la base de cálculo de rotación del equipo clave. El texto del compromiso y la nómina de la Tabla~\ref{tab:12D-compromisos} constituyen la carta de compromiso del equipo, que acompaña al Formulario T-8 junto con los currículos del Anexo 12.B (Bases Administrativas, art.~34.1, p.~22). Nombre, rol, dedicación y período de cada integrante coinciden con su nominación en el T-8.

\section*{Referencias}
\phantomsection
\addcontentsline{toc}{section}{Referencias}

Las condiciones contractuales y los requisitos técnicos citados corresponden a los siguientes documentos de la licitación:

\begin{itemize}
\item Distribuidora Puelche S.A. (2026a). \emph{Bases Administrativas}.
\item Distribuidora Puelche S.A. (2026b). \emph{Bases Técnicas Transversales}.
\end{itemize}

\section*{Declaración de uso de IA}
\phantomsection
\addcontentsline{toc}{section}{Declaración de uso de IA}

La declaración consolidada al final del Subdocumento 12 identifica, para los Anexos 12.A a 12.D, las herramientas utilizadas, la finalidad de su uso y el alcance de la revisión humana.
